\chapter{Metodologia}\label{cap:metodologia}

Neste capítulo, descreveu-se o processamento executado sobre os panoramas de escoamento
bifásico do LETEF, com os resultados preservados em
\codigo{src2/v16\_spatial2\_18\_71}. Aplicou-se a formulação alternativa de
\citeonline{schuck2019} a sinais espaciais de intensidade e compararam-se dois
indicadores: o escore de coerência entre pares de ensaios e o coeficiente de Pearson
entre descritores espectrais individuais. A execução foi exploratória e descritiva,
sem treinamento de classificador ou teste de significância estatística.

% =========================================================
\section{Visão geral do método executado}\label{sec:met-visao}
% =========================================================

O processamento partiu dos panoramas existentes e releu seus pixels para extrair a
média das linhas 18 a 71, inclusive. No \autoref{qua:met-etapas}, distinguiram-se as duas
ramificações da análise. Essa distinção foi necessária porque o espectro individual
foi calculado uma única vez por ensaio completo, enquanto os mapas de coerência
usaram o comprimento comum de cada par.

\begin{quadro}[htb]
  \centering
  \caption{Etapas do processamento efetivamente executado.}
  \label{qua:met-etapas}
  \begin{tabular}{p{3.1cm}p{10.5cm}}
    \hline
    \textbf{Etapa} & \textbf{Operação e finalidade} \\
    \hline
    Entrada & Conferência dos panoramas e extração da média de 54 linhas por coluna. \\
    Preparação & Inversão da intensidade, remoção de tendência linear, centralização e padronização. \\
    Perfil individual & Autocorrelação e CWT do panorama completo; perfil de potência no atraso nulo para o descritor de Pearson. \\
    Comparação de pares & Corte no menor comprimento, preparação dos sinais, autocorrelações, correlação cruzada e campos WPS, WCS e WCO. \\
    Resumos & Coeficiente de Pearson dos resíduos espectrais e escore dos máximos de WCO ponderados pela potência. \\
    Avaliação & Comparações dentro e entre padrões, médias por ensaio e controles de sensibilidade e de qualidade. \\
    \hline
  \end{tabular}
  \fonte{Elaborado pelo autor a partir dos scripts e da configuração da execução.}
\end{quadro}

% =========================================================
\section{Dados experimentais e seleção dos ensaios}\label{sec:met-dados}
% =========================================================

Utilizaram-se sequências de imagens do arranjo de isobutano R600a em tubo horizontal
de 1\,mm, descrito por \citeonline{deoliveira2021} e apresentado na
\autoref{sec:arranjo}. A auditoria local compreendeu 31 registros, cada um com 501
quadros de $1016\times80$ pixels. A seleção preservada na execução reuniu 26 ensaios:
dez de \textit{churn}, nove de anular ondulado, três de \textit{plug}, três de
\textit{slug} e um de bolha dispersa. Os rótulos dos 26 registros concordaram com a
tabela manual local; não foram inferidos dos valores de coerência ou de correlação.

Permaneceram fora das estatísticas os registros \codigo{m03q01\_001},
\codigo{m03q80\_001}, \codigo{m04q05\_001}, \codigo{m05q05\_002} e
\codigo{m07q10\_001}. A auditoria preservou sua identificação e constatou, no
registro \codigo{m05q05\_002}, divergência entre o rótulo local de bolha dispersa e
a classificação \textit{slug} da condição na tabela manual. Essa divergência não
foi usada para reclassificar ou reincluir o registro.

Formaram-se todas as combinações de dois ensaios distintos: $26\times25/2=325$
pares únicos. Destes, 87 foram intrapadrão, sendo 45 de \textit{churn}, 36 de anular
ondulado, três de \textit{plug} e três de \textit{slug}; os outros 238 foram
interpadrões. O único ensaio de bolha dispersa participou de 25 comparações
interpadrões, mas não permitiu avaliar repetibilidade dentro desse padrão.

Os metadados das aquisições incluídas registraram taxas de 4\,000, 5\,000 e
7\,000 quadros/s e tempos de exposição de 62, 127 e 132\,$\mu$s. Essas informações
foram conferidas nos dois arquivos de configuração de cada aquisição. Não se
aplicou uma correção fotométrica baseada nesses códigos, nem se converteu o eixo
espacial em frequência temporal a partir da taxa de quadros.

% ---------------------------------------------------------
\subsection{Reconstrução e extração do sinal espacial}\label{sec:met-extracao}
% ---------------------------------------------------------

Os panoramas utilizados já haviam sido montados a partir de tiras centrais dos
quadros. O algoritmo estimou deslocamentos entre quadros consecutivos por
correlação dos perfis médios das linhas 20 a 59 e utilizou esses deslocamentos como
larguras das tiras. Larguras não positivas foram substituídas por 1 pixel; na última
tira, repetiu-se a largura estimada no último intervalo. A junção incluiu a mistura
linear de até dois pixels nas bordas das tiras. Na auditoria posterior, essa
montagem foi reproduzida em memória e comparada pixel a pixel com os 31 panoramas
existentes, sem substituir as entradas da análise.

Para a análise espectral desta execução, a faixa vertical foi ampliada para as
linhas 18 a 71, inclusive, com índices iniciados em zero. O sinal bruto do ensaio
$i$ foi definido pela média aritmética:

\begin{equation}\label{eq:met-extracao}
  I_i[n]=\frac{1}{54}\sum_{v=18}^{71} A_i[v,n],
  \qquad n=0,\ldots,N_i-1,
\end{equation}

\noindent
em que $A_i[v,n]$ representou a intensidade do pixel da linha $v$ e coluna $n$ do panorama,
$I_i[n]$ representou a intensidade média dessa coluna e $N_i$, sua largura. As planilhas
anteriores, correspondentes às linhas 20 a 59, serviram apenas à conferência de
procedência; não forneceram o sinal da \autoref{eq:met-extracao}.

As larguras dos 26 panoramas variaram de 3\,171 a 23\,079 pixels. O passo espacial
adotado foi de 1 pixel do panorama reconstruído. Esse eixo reuniu observações de
instantes diferentes, posicionadas pelo deslocamento óptico estimado; portanto,
não correspondeu a uma fotografia espacial instantânea nem a uma medida calibrada
de distância física. O maior número de colunas também não representou o mesmo
número de observações independentes. A intensidade média permaneceu uma grandeza
óptica, sem conversão direta em fração de vazio.

% =========================================================
\section{Preparação dos sinais}\label{sec:met-preproc}
% =========================================================

Para o perfil individual, preparou-se o sinal completo de cada ensaio. Para os
mapas de um par $x,y$, selecionaram-se primeiro os prefixos de comprimento
$N=\min(N_x,N_y)$, e somente depois se preparou cada prefixo. Não houve alinhamento
de eventos ou escolha de trechos em função da semelhança obtida.

A preparação consistiu nas operações:

\begin{align}
  J[n]&=255-I[n], \label{eq:met-inversao}\\
  d[n]&=J[n]-(\widehat\alpha+\widehat\beta n), \label{eq:met-tendencia}\\
  x[n]&=\frac{d[n]-\bar d}{s_d},\qquad
  s_d=\sqrt{\frac{1}{N}\sum_{n=0}^{N-1}(d[n]-\bar d)^2}.
  \label{eq:met-padronizacao}
\end{align}

\noindent
A intensidade invertida $J[n]$ associou valores maiores aos pixels mais escuros.
Os coeficientes $\widehat\alpha$ e $\widehat\beta$ definiram a reta ajustada por
mínimos quadrados, isto é, a reta que minimizou a soma dos desvios quadráticos.
O sinal sem tendência $d[n]$ foi centralizado pela média $\bar d$ e dividido pelo
desvio padrão populacional $s_d$, calculado com divisor $N$. O sinal preparado
$x[n]$ apresentou média nula e variância unitária, a menos de arredondamento
numérico. Aplicou-se a mesma operação ao segundo sinal $y[n]$.

Esse preparo reduziu diferenças de nível e amplitude e removeu a tendência linear
no trecho analisado. Não comprovou estacionariedade, não recuperou informação
perdida por saturação e não corrigiu os deslocamentos da montagem. A formulação
baseada em funções de correlação pressupôs processos estacionários
\cite{schuck2019}; no domínio reconstruído, a adequação dessa hipótese permaneceu
uma condição de interpretação, e não uma consequência automática da padronização.

% =========================================================
\section{Correlações e transformada \textit{Wavelet}}\label{sec:met-wco}
% =========================================================

Para sinais de comprimento comum $N$, estimaram-se as correlações bilaterais:

\begin{equation}\label{eq:crosscorr-disc}
  \widehat R_{xy}[m]=\frac{1}{N}
  \sum_{n=\max(0,-m)}^{\min(N-1,N-1-m)}
  x[n+m]y^*[n],\qquad -(N-1)\leq m\leq N-1.
\end{equation}

\begin{equation}\label{eq:autocorr-disc}
  \widehat R_{xx}[m]=\left.\widehat R_{xy}[m]\right|_{y=x},
  \qquad
  \widehat R_{yy}[m]=\left.\widehat R_{xy}[m]\right|_{x=y}.
\end{equation}

\noindent
O atraso espacial $m$ foi expresso em pixels. Os limites da soma garantiram que
ambos os índices estivessem dentro do registro, inclusive para atrasos negativos.
O asterisco indicou conjugação complexa, sem efeito sobre os sinais reais de
intensidade. Nessa convenção, quando $y$ foi uma cópia de $x$ deslocada para a
direita por $d$ amostras, o pico ocorreu em atraso espacial $m=-d$.

Empregou-se a normalização enviesada, com divisor $N$ em todos os atrasos. Essa
escolha diferiu dos estimadores com divisor $N-m$, apresentados para atrasos
não negativos nas Equações 32.21 a 32.23 de \citeonline{schuck2019}.
Na extensão para atrasos positivos e negativos, esse divisor corresponderia a
$N-|m|$. Tratou-se, portanto, de uma adaptação declarada da
estimação das correlações, mantendo-se as definições alternativas de WPS, WCS e WCO.
Nos atrasos extremos, a soma envolveu poucos produtos, mas permaneceu dividida pelo
comprimento total do registro.

% ---------------------------------------------------------
\subsection{Parametrização e implementação da Morlet}\label{sec:met-wco-morlet}
% ---------------------------------------------------------

Utilizou-se a Morlet complexa da \autoref{eq:morlet}, com frequência angular central
adimensional $\omega_0=6$. A grade compreendeu 200 comprimentos de onda
$\lambda_k$, igualmente espaçados em escala logarítmica entre 50 e 3\,000 pixels.
A conversão entre a escala $a_k$, expressa em amostras, e o comprimento de onda foi
\cite{torrence1998,schuck2019}:

\begin{equation}\label{eq:met-escala}
  \lambda_k=\frac{4\pi}{6+\sqrt{38}}\,a_k\,\Delta s,
  \qquad \Delta s=1\,\mathrm{px}.
\end{equation}

\noindent
O fator $\Delta s$ representou o espaçamento de amostragem. Assim, a escala e o
comprimento de onda não foram tratados como valores numericamente idênticos; cada
escala correspondeu a aproximadamente $1{,}03a_k$ pixels.

A CWT foi implementada no domínio de Fourier, com \textit{NumPy} e
\textit{SciPy}, pelo multiplicador da Morlet:

\begin{equation}\label{eq:met-morlet-fourier}
  H_{a_k}(\omega)=\sqrt{2\pi a_k}\,\pi^{-1/4}
       \exp\!\left[-\frac{(a_k\omega-6)^2}{2}\right].
\end{equation}

\noindent
A frequência angular digital $\omega$ foi expressa em radianos por amostra.
Para cada escala, multiplicou-se a transformada de Fourier da correlação pelo
filtro $H_{a_k}$ e aplicou-se a transformada inversa. O comprimento da FFT foi
estendido por $2\lceil8a_{\max}\rceil$ zeros, antes do arredondamento para um
tamanho computacionalmente eficiente, em que $a_{\max}$ representou a maior
escala e $\lceil\cdot\rceil$, o arredondamento para o inteiro superior.
Essa implementação realizou a Morlet da \autoref{eq:morlet}.

Como expressão discreta de referência, para uma função de correlação
$\widehat R_{xy}$, utilizou-se:

\begin{equation}\label{eq:cwt-disc}
  W_{R_{xy}}(a_k,b)\approx\frac{1}{\sqrt{a_k}}
     \sum_{m=-(N-1)}^{N-1}\widehat R_{xy}[m]\,
     \psi^*\!\left(\frac{m-b}{a_k}\right).
\end{equation}

\noindent
A coordenada $b$ indicou a posição no eixo de atrasos da correlação, também em
amostras espaciais, e não a posição original no panorama. A soma direta com a
Morlet serviu à verificação independente da implementação espectral.

% ---------------------------------------------------------
\subsection{Campos e perfis espectrais}\label{sec:met-analise}
% ---------------------------------------------------------

Aplicaram-se as \autoref{eq:wps-alt}, \autoref{eq:wcs-alt} e
\autoref{eq:wco-alt} às correlações estimadas. O campo de potência alternativa
$\widetilde P_x(a,b)$ foi o módulo da CWT da autocorrelação, sem elevar esse módulo
novamente ao quadrado. O campo cruzado $\widetilde C_{xy}(a,b)$ permaneceu complexo;
seu módulo foi utilizado nas figuras. O campo de coerência
$\widetilde\gamma_{xy}(a,b)$ variou com a escala e o atraso, enquanto seu denominador
utilizou as duas potências no atraso nulo.

Os mapas compreenderam todos os atrasos de $-(N-1)$ a $N-1$, sem máscara
geométrica. Para proteger a divisão, exigiu-se que a potência no atraso nulo de
cada sinal superasse $10^{-10}$ do respectivo máximo. Essa verificação numérica
não excluiu nenhuma das 200 escalas nos 325 pares. Os valores de coerência não
foram truncados artificialmente em 1.

Distinguiram-se dois perfis individuais:

\begin{equation}\label{eq:met-perfis}
  P_{ik}=\widetilde P_i(a_k,0),\qquad
  M_{ik}=\frac{1}{2N_i-1}\sum_{b=-(N_i-1)}^{N_i-1}\widetilde P_i(a_k,b).
\end{equation}

\noindent
A potência no atraso nulo $P_{ik}$ forneceu a entrada do descritor de Pearson.
A potência média nos atrasos $M_{ik}$ resumiu o campo em uma curva, exibida ao lado
dos mapas. Como essas operações foram diferentes, seus picos não precisaram
coincidir. O comprimento de onda do pico foi obtido pelo argumento do máximo na
grade utilizada, especificando-se em cada caso qual dos dois perfis foi analisado.

% =========================================================
\section{Escore de coerência entre ensaios}\label{sec:met-escore}
% =========================================================

Para resumir cada campo de coerência, calculou-se o máximo em atraso para cada
escala, $c_k=\max_b\widetilde\gamma_{xy}(a_k,b)$. A importância da escala foi
ponderada pela média geométrica das potências no atraso nulo,
$w_k=\sqrt{\widetilde P_x(a_k,0)\widetilde P_y(a_k,0)}$. Selecionaram-se as
escalas em que ambos os sinais apresentaram potência suficiente:

\begin{equation}\label{eq:met-piso}
  \mathcal K_\eta=\left\{k:
  \widetilde P_x(a_k,0)\geq\eta\max_j\widetilde P_x(a_j,0),\quad
  \widetilde P_y(a_k,0)\geq\eta\max_j\widetilde P_y(a_j,0)\right\}.
\end{equation}

\begin{equation}\label{eq:met-escore}
  S_{xy}^{(\eta)}=
  \frac{\sum_{k\in\mathcal K_\eta}w_kc_k}
       {\sum_{k\in\mathcal K_\eta}w_k}.
\end{equation}

\noindent
O piso relativo de potência $\eta$ foi 0,05 no escore principal de coerência
$S_{xy}^{(0,05)}$. Também foram registrados os escores com pisos de 0,01 e 0,10.
O peso $w_k$ reduziu a contribuição de escalas de pequena potência, e o conjunto
$\mathcal K_\eta$ exigiu participação de ambos os sinais. O piso relativo de
potência não foi um limiar de significância e não removeu regiões dos mapas.
O escore foi um resumo adotado nesta execução, não uma equação adicional atribuída
a \citeonline{schuck2019}.

Cada máximo por escala pôde ocorrer em um atraso diferente. Assim, um escore de
coerência elevado não correspondeu necessariamente a um único alinhamento entre os
sinais. Registraram-se ainda o máximo integral, a média integral e a média
ponderada da coerência no atraso nulo como medidas complementares; nenhuma delas
substituiu o escore principal de coerência.

% =========================================================
\section{Coeficiente de Pearson dos descritores espectrais}\label{sec:met-pearson}
% =========================================================

Para cada ensaio completo, transformou-se a potência no atraso nulo $P_{ik}$ em
logaritmo natural e removeu-se uma reta em função do logaritmo do comprimento de
onda. Os coeficientes da reta foram calculados pela minimização:

\begin{equation}\label{eq:met-ajuste}
  (\widehat c_i,\widehat d_i)=\underset{c,d}{\arg\min}
  \sum_{k=1}^{200}\left[\ln(P_{ik}+\varepsilon)-c-d\ln(\lambda_k)\right]^2,
  \qquad \varepsilon=10^{-12}.
\end{equation}

\begin{equation}\label{eq:met-residuos}
  z_{ik}=\ln(P_{ik}+\varepsilon)-\widehat c_i-\widehat d_i\ln(\lambda_k).
\end{equation}

\noindent
O termo $\varepsilon$ evitou o logaritmo de zero. O coeficiente $\widehat c_i$
representou o nível ajustado e o coeficiente $\widehat d_i$, a inclinação ajustada.
O resíduo espectral $z_{ik}$ expressou quanto cada escala ficou acima ou abaixo
dessa tendência. A operação retirou o nível e a inclinação globais do perfil,
concentrando a comparação nas variações remanescentes. Os comprimentos de onda
foram inseridos numericamente em pixels, com a mesma unidade em todos os ensaios.

Entre os resíduos de dois ensaios $i,j$, calculou-se o coeficiente de Pearson
$r_{ij}$, conforme sua definição centrada \cite{cunningham2020}:

\begin{equation}\label{eq:met-pearson}
  r_{ij}=\frac{\sum_{k=1}^{200}(z_{ik}-\bar z_i)(z_{jk}-\bar z_j)}
   {\sqrt{\sum_{k=1}^{200}(z_{ik}-\bar z_i)^2}
    \sqrt{\sum_{k=1}^{200}(z_{jk}-\bar z_j)^2}}.
\end{equation}

\noindent
As médias $\bar z_i$ e $\bar z_j$ foram calculadas nas 200 escalas. O coeficiente
de Pearson $r_{ij}$ próximo de 1 indicou variações residuais no mesmo sentido;
próximo de $-1$, variações em sentidos opostos; e próximo de zero, pequena
associação linear entre os resíduos \cite{cunningham2020}. Preservou-se o sinal
do coeficiente de Pearson, sem tomar seu módulo.

Cada descritor foi mantido fixo para todos os parceiros do ensaio. Não se calculou
o coeficiente de Pearson diretamente entre intensidades, campos de coerência ou
perfis de potência média nos atrasos. A semelhança medida referiu-se ao descritor
residual e não conservou a amplitude média, a inclinação removida ou a fase dos
sinais originais.

% =========================================================
\section{Agregações e controles de sensibilidade}\label{sec:met-agregacao}
% =========================================================

Calcularam-se médias e medianas dentro e entre padrões, além das médias de cada
ensaio com seus parceiros. A diagonal unitária da matriz do coeficiente de
Pearson foi excluída dessas médias. A média geral de cada ensaio reuniu seus 25
parceiros; a média intrapadrão de bolha dispersa permaneceu indefinida.

Foram utilizadas duas ponderações. A média por par atribuiu peso igual a cada
comparação. A média balanceada calculou primeiro a média de cada grupo e depois
atribuiu peso igual aos grupos: quatro grupos intrapadrão e dez combinações
interpadrões. Para uma métrica $h$, representando o coeficiente de Pearson ou o
escore de coerência, definiu-se o contraste descritivo:

\begin{equation}\label{eq:met-contraste}
  \Delta_h=\overline h_{\mathrm{intra}}-\overline h_{\mathrm{inter}}.
\end{equation}

\noindent
O contraste descritivo $\Delta_h$ positivo indicou maior média dentro dos padrões;
o contraste negativo indicou maior média entre padrões. As duas ponderações
responderam a perguntas distintas: uma descreveu o conjunto de pares, e a outra
reduziu o predomínio dos grupos mais numerosos. Nenhuma forneceu, por si só,
uma medida de significância estatística ou de acurácia.

A auditoria de qualidade motivou controles explícitos para o registro anular
\codigo{m03q40\_001}. Compararam-se, para os nove ensaios anulares, o prefixo comum
de 14\,008 pixels anterior à primeira substituição de deslocamento, a faixa de
comprimentos de onda de 50 a 1\,024 pixels, a preparação sem remoção de tendência
linear, a extração pelas linhas 20 a 59 e a equalização global do histograma.
Cada controle aplicou a mesma regra a todos os nove ensaios. Na equalização,
a transformação foi estimada separadamente em cada aquisição a partir do
histograma da faixa 18 a 71; ela permaneceu um controle exploratório.

Também se recalcularam as agregações sem \codigo{m03q40\_001}, preservando os
descritores dos outros ensaios. A referência principal permaneceu com 26 ensaios;
a variante com 25 ensaios não foi escolhida para maximizar a separação. Esses
controles avaliaram sensibilidade e não estabeleceram uma causa única para o
comportamento do registro.

% =========================================================
\section{Verificação e alcance da avaliação}\label{sec:met-verificacao}
% =========================================================

Confrontaram-se as correlações com somas diretas e a CWT espectral com integração
temporal independente da Morlet. Verificaram-se a autocoerência no atraso nulo,
a recuperação de um atraso conhecido e a reflexão do mapa ao trocar a ordem dos
sinais. Os 325 coeficientes de Pearson foram comparados com uma rotina
independente da biblioteca \textit{SciPy}; as matrizes e as agregações também
foram verificadas numericamente.

Os resumos integrais foram calculados antes da redução dos mapas para
apresentação. Os arquivos usados para reproduzir as figuras conservaram cerca de
1\,800 posições de atraso, incluindo o zero e os extremos. Por isso, os máximos
integrais não foram recalculados apenas dos mapas reduzidos.

A faixa de 50 a 3\,000 pixels foi mantida inclusive nos registros mais curtos.
O menor panorama conteve aproximadamente 1,057 comprimentos de onda de
3\,000 pixels, o que limitou o suporte às maiores escalas. A extensão bilateral
da correlação não acrescentou novas observações. Além disso, os pares
compartilharam ensaios e as escalas da CWT foram correlacionadas. Por essas
razões, a avaliação relatou tendências descritivas, sem tratar os 325 pares ou
as 200 escalas como réplicas independentes para calcular valores de significância
ou intervalos de confiança. Esse cuidado foi apoiado no princípio de evitar
pseudorreplicação, isto é, tratar medidas sem independência como réplicas
independentes em testes inferenciais \cite{hurlbert1984}.
Os resultados constituíram uma base para análises
adicionais de qualidade, robustez e discriminação.
